\section{Diario della riunione}

\subsection{Dubbi sul corretto utilizzo dei database vettoriali}
Abbiamo chiesto come inizializzare i DB vettoriali, in quanto il nostro metodo di lavoro è indicizzare e trasformare in vettori. Il gruppo chiede se il metodo è giusto oppure ci sono vie più semplici.\\
Ci è stato detto che vie più semplici purtroppo non esistono e, mostrando il codice, ci è stata confermata la correttezza del nostro operato.\\
Un possibile metodo sarebbe da utilizzare solo il DB relazionale, cosa che è stata prontamente scartata dal gruppo in quanto le performance risultano molto più lente.\\
Successivamente abbiamo chiesto che cosa ci va nei DB vettoriali, ci è stato suggerito di inserire sicuramente tutte le cose relative agli articoli, e in generale tutto ciò che è utile allo scopo di costruzione dell'ordine.\\
Infatti, tutto ciò che fa parte dei prodotti va sicuramente inserito.\\
Sarà utile pure inserire lo storico degli ordini, di sicuro non vanno inserite cose relative agli utenti come la password e la mail.

\subsection{Considerazione delle unità di misura per l'MVP}
Abbiamo chiesto, siccome nel dataset fornitoci sono stati inserite le unità di misure di ogni prodotto, se queste fossero necessarie da inserire per l'approvazione dell'MVP.\\
La risposta è stata che non è una cosa necessaria ai fini dell'MVP, e che se ci avanzerà tempo potremmo inserirla in quanto preferibile.

\subsection{Tipo di architettura di deployment}
Abbiamo chiesto se fosse preferibile per il capitolato avere un'architettura monolitica oppure un'architettura a microservizi.\\
Ci è stato risposto che sarebbe meglio il monolite, non perché sia più facile ma perché più utile al nostro scopo siccome il prodotto non ha numerose funzionalità che necessitano di microservizi e non c'è bisogno di scalare per il futuro.

\subsection{Suggerimenti per l'architettura}
Abbiamo mostrato tutte le architetture fatte fino a quel momento, mostrando i diagrammi delle classi comprensivi dei design pattern.\\
Ci è stato detto che tutte le nostre idee sono buone e valide e che come inizio non va bene.\\
Poi abbiamo chiesto se è possibile, basandosi sulla loro esperienza, se si può riconoscere se una progettazione è sbagliata prima dell'effettiva codifica. \\
Ci è stato detto che è quasi impossibile a colpo d'occhio a meno che non siano errori grossolani, quindi uno dei modi è progettare è tentare.